\textbf{Phần 1. Trắc nghiệm lựa chọn}
\setcounterpageref{ex}{0}

\begin{ex} %[5D10N57]
Kết quả của phép tính $3$ giờ $15$ phút $+ 2$ giờ $20$ phút là
\choice
{$5$ giờ $25$ phút}
{\True $5$ giờ $35$ phút}
{$5$ giờ $45$ phút}
{$6$ giờ $05$ phút}
\loigiai{
$3$ giờ $15$ phút $+ 2$ giờ $20$ phút $= 5$ giờ $35$ phút.\\
\textbf{Đáp án của câu hỏi là B}
}
\end{ex}

\begin{ex} %[5D10N57]
Kết quả của phép tính $8$ phút $45$ giây $- 3$ phút $20$ giây là
\choice
{\True $5$ phút $25$ giây}
{$5$ phút $15$ giây}
{$5$ phút $35$ giây}
{$6$ phút $25$ giây}
\loigiai{
$8$ phút $45$ giây $- 3$ phút $20$ giây $= 5$ phút $25$ giây.\\
\textbf{Đáp án của câu hỏi là A}
}
\end{ex}

\begin{ex} %[5D10H57]
Kết quả của phép tính $4$ giờ $40$ phút $+ 3$ giờ $35$ phút sau khi đổi đơn vị là
\choice
{$7$ giờ $75$ phút}
{$8$ giờ $25$ phút}
{$7$ giờ $15$ phút}
{\True $8$ giờ $15$ phút}
\loigiai{
$4$ giờ $40$ phút $+ 3$ giờ $35$ phút $= 7$ giờ $75$ phút $= 8$ giờ $15$ phút.\\
\textbf{Đáp án của câu hỏi là D}
}
\end{ex}

\begin{ex} %[5D10H57]
Thực hiện phép trừ: $6$ ngày $8$ giờ $- 2$ ngày $15$ giờ. Kết quả là
\choice
{$4$ ngày $7$ giờ}
{\True $3$ ngày $17$ giờ}
{$3$ ngày $7$ giờ}
{$4$ ngày $17$ giờ}
\loigiai{
$6$ ngày $8$ giờ $= 5$ ngày $32$ giờ.\\
$5$ ngày $32$ giờ $- 2$ ngày $15$ giờ $= 3$ ngày $17$ giờ.\\
\textbf{Đáp án của câu hỏi là B}
}
\end{ex}

\begin{ex} %[5D10H57]
Một xe ô tô xuất phát lúc $7$ giờ $15$ phút và đến nơi lúc $10$ giờ $45$ phút. Thời gian ô tô đi trên đường là
\choice
{$3$ giờ}
{$3$ giờ $45$ phút}
{$4$ giờ}
{\True $3$ giờ $30$ phút}
\loigiai{
$10$ giờ $45$ phút $- 7$ giờ $15$ phút $= 3$ giờ $30$ phút.\\
\textbf{Đáp án của câu hỏi là D}
}
\end{ex}

\begin{ex} %[5D10H57]
Bạn Lan học bài từ $19$ giờ $30$ phút đến $21$ giờ $15$ phút. Thời gian Lan học bài là
\choice
{\True $1$ giờ $45$ phút}
{$1$ giờ $15$ phút}
{$2$ giờ}
{$2$ giờ $15$ phút}
\loigiai{
$21$ giờ $15$ phút $- 19$ giờ $30$ phút $= 20$ giờ $75$ phút $- 19$ giờ $30$ phút $= 1$ giờ $45$ phút.\\
\textbf{Đáp án của câu hỏi là A}
}
\end{ex}

\begin{ex} %[5D10V57]
Một người đi từ A lúc $8$ giờ $20$ phút, đi hết $2$ giờ $45$ phút thì đến B. Người đó đến B lúc mấy giờ?
\choice
{$10$ giờ $5$ phút}
{$11$ giờ $15$ phút}
{\True $11$ giờ $5$ phút}
{$10$ giờ $45$ phút}
\loigiai{
$8$ giờ $20$ phút $+ 2$ giờ $45$ phút $= 10$ giờ $65$ phút $= 11$ giờ $5$ phút.\\
\textbf{Đáp án của câu hỏi là C}
}
\end{ex}

\begin{ex} %[5D10V57]
Bác Thợ may chiếc áo hết $1$ giờ $45$ phút, may chiếc quần hết $1$ giờ $20$ phút. Tổng thời gian bác may cả áo và quần là
\choice
{\True $3$ giờ $5$ phút}
{$2$ giờ $5$ phút}
{$3$ giờ $15$ phút}
{$2$ giờ $55$ phút}
\loigiai{
$1$ giờ $45$ phút $+ 1$ giờ $20$ phút $= 2$ giờ $65$ phút $= 3$ giờ $5$ phút.\\
\textbf{Đáp án của câu hỏi là A}
}
\end{ex}

\begin{ex} %[5D10V57]
Trong một hội thao chạy tiếp sức $4 \times 100\text{ m}$, vận động viên thứ nhất chạy hết $14{,}2$ giây, vận động viên thứ hai chạy hết $13{,}8$ giây, vận động viên thứ ba chạy hết $14{,}5$ giây, vận động viên thứ tư chạy hết $13{,}5$ giây. Tổng thời gian chạy của cả đội là
\choice
{$55$ giây}
{$57$ giây}
{\True $56$ giây}
{$58$ giây}
\loigiai{
$14{,}2 + 13{,}8 + 14{,}5 + 13{,}5 = 56$ giây.\\
\textbf{Đáp án của câu hỏi là C}
}
\end{ex}

\begin{ex} %[5D10V57]
Một xe buýt bắt đầu xuất phát từ bến A lúc $6$ giờ sáng và cứ sau $45$ phút lại hoàn thành một chuyến đi về bến A. Chuyến xe thứ tư trở về bến A lúc mấy giờ?
\choice
{$8$ giờ $30$ phút}
{\True $9$ giờ}
{$9$ giờ $15$ phút}
{$9$ giờ $30$ phút}
\loigiai{
Thời gian đi 4 chuyến: $45 \times 4 = 180$ phút $= 3$ giờ.\\
Thời điểm về bến sau 4 chuyến: $6$ giờ $+ 3$ giờ $= 9$ giờ.\\
\textbf{Đáp án của câu hỏi là B}
}
\end{ex}

\begin{ex} %[5D10V57]
Tính giá trị biểu thức: $12$ giờ $30$ phút $- (4$ giờ $15$ phút $+ 3$ giờ $40$ phút$)$.
\choice
{$5$ giờ $5$ phút}
{$4$ giờ $45$ phút}
{\True $4$ giờ $35$ phút}
{$5$ giờ $15$ phút}
\loigiai{
$4$ giờ $15$ phút $+ 3$ giờ $40$ phút $= 7$ giờ $55$ phút.\\
$12$ giờ $30$ phút $- 7$ giờ $55$ phút $= 11$ giờ $90$ phút $- 7$ giờ $55$ phút $= 4$ giờ $35$ phút.\\
\textbf{Đáp án của câu hỏi là C}
}
\end{ex}

\begin{ex} %[5D10C57]
Một chuyến tàu đi từ Hà Nội đến Đà Nẵng xuất phát lúc $19$ giờ $30$ phút ngày hôm trước và đến nơi lúc $11$ giờ $15$ phút trưa ngày hôm sau. Tổng thời gian chuyến đi của tàu là
\choice
{$14$ giờ $45$ phút}
{$16$ giờ $15$ phút}
{$15$ giờ $15$ phút}
{\True $15$ giờ $45$ phút}
\loigiai{
Thời gian đi từ $19$ giờ $30$ phút đến $24$ giờ: $24$ giờ $- 19$ giờ $30$ phút $= 4$ giờ $30$ phút.\\
Thời gian đi ngày hôm sau: $11$ giờ $15$ phút.\\
Tổng thời gian: $4$ giờ $30$ phút $+ 11$ giờ $15$ phút $= 15$ giờ $45$ phút.\\
\textbf{Đáp án của câu hỏi là D}
}
\end{ex}

\textbf{Phần 2. Trắc nghiệm đúng sai}
\setcounterpageref{ex}{0}

\begin{ex} %[5D10H57]
Xét tính đúng sai của các phép tính thời gian:
\choiceTF
{\True $5$ giờ $35$ phút $+ 3$ giờ $45$ phút $= 9$ giờ $20$ phút.}
{$7$ phút $15$ giây $- 4$ phút $45$ giây $= 3$ phút $30$ giây.}
{\True $7$ phút $15$ giây $- 4$ phút $45$ giây $= 2$ phút $30$ giây.}
{\True $8$ ngày $6$ giờ $- 3$ ngày $18$ giờ $= 4$ ngày $12$ giờ.}
\loigiai{
\begin{itemize}
    \item $5$h $35$p $+ 3$h $45$p $= 8$h $80$p $= 9$h $20$p. Đúng.
    \item $7$p $15$s $- 4$p $45$s $= 6$p $75$s $- 4$p $45$s $= 2$p $30$s $\neq 3$p $30$s. Ý 2 Sai, Ý 3 Đúng.
    \item $8$ ngày $6$h $- 3$ ngày $18$h $= 7$ ngày $30$h $- 3$ ngày $18$h $= 4$ ngày $12$h. Đúng.
\end{itemize}
\begin{itemize}
    \item a) \textbf{Đúng}.
    \item b) \textbf{Sai}.
    \item c) \textbf{Đúng}.
    \item d) \textbf{Đúng}.
\end{itemize}
}
\end{ex}

\begin{ex} %[5D10V57]
An bắt đầu làm bài kiểm tra lúc $8$ giờ $10$ phút và nộp bài lúc $9$ giờ $05$ phút. Bình làm bài từ $8$ giờ $15$ phút và nộp bài lúc $9$ giờ $00$ phút. Xét tính đúng sai:
\choiceTF
{\True Thời gian An làm bài là $55$ phút.}
{\True Thời gian Bình làm bài là $45$ phút.}
{An làm bài nhanh hơn Bình $10$ phút.}
{\True Bình làm bài nhanh hơn An $10$ phút.}
\loigiai{
\begin{itemize}
    \item Thời gian An làm bài: $9$h $05$p $- 8$h $10$p $= 55$ phút. Đúng.
    \item Thời gian Bình làm bài: $9$h $00$p $- 8$h $15$p $= 45$ phút. Đúng.
    \item Bình làm mất ít thời gian hơn nên Bình làm nhanh hơn An $55 - 45 = 10$ phút. Ý 3 Sai, Ý 4 Đúng.
\end{itemize}
\begin{itemize}
    \item a) \textbf{Đúng}.
    \item b) \textbf{Đúng}.
    \item c) \textbf{Sai}.
    \item d) \textbf{Đúng}.
\end{itemize}
}
\end{ex}

\textbf{Phần 3. Trả lời ngắn}
\setcounterpageref{ex}{0}

\begin{ex} %[5D10H57]
Tính tổng: $4$ giờ $45$ phút $+ 2$ giờ $35$ phút theo phút.
\shortans{440}
\loigiai{
\textbf{Đáp số:} $440$\\
$7$ giờ $20$ phút $= 7 \times 60 + 20 = 440$ phút.
}
\end{ex}

\begin{ex} %[5D10H57]
Tính hiệu: $9$ giờ $15$ phút $- 5$ giờ $40$ phút theo phút.
\shortans{215}
\loigiai{
\textbf{Đáp số:} $215$\\
$8$ giờ $75$ phút $- 5$ giờ $40$ phút $= 3$ giờ $35$ phút $= 215$ phút.
}
\end{ex}

\begin{ex} %[5D10V57]
Một cuộc họp dự kiến bắt đầu lúc $14$ giờ $15$ phút và kéo dài trong $1$ giờ $45$ phút. Hỏi cuộc họp dự kiến kết thúc lúc mấy giờ? (Chỉ ghi số giờ, ví dụ: 16).
\shortans{16}
\loigiai{
\textbf{Đáp số:} $16$\\
$14$ giờ $15$ phút $+ 1$ giờ $45$ phút $= 16$ giờ.
}
\end{ex}

\begin{ex} %[5D10V57]
Thời gian xem một bộ phim là $1$ giờ $50$ phút. Bộ phim kết thúc lúc $21$ giờ $30$ phút. Hỏi bộ phim bắt đầu chiếu lúc mấy giờ mấy phút? Tính số phút từ $0$ giờ đến lúc bắt đầu chiếu.
\shortans{1180}
\loigiai{
\textbf{Đáp số:} $1180$\\
Bắt đầu chiếu: $21$ giờ $30$ phút $- 1$ giờ $50$ phút $= 19$ giờ $40$ phút $= 19 \times 60 + 40 = 1\,180$ phút.
}
\end{ex}

\begin{ex} %[5D10V57]
Một công nhân làm việc ca 1 hết $3$ giờ $45$ phút, ca 2 hết $4$ giờ $15$ phút. Hỏi cả hai ca công nhân đó làm việc hết bao nhiêu giờ?
\shortans{8}
\loigiai{
\textbf{Đáp số:} $8$\\
$3$ giờ $45$ phút $+ 4$ giờ $15$ phút $= 8$ giờ.
}
\end{ex}

\begin{ex} %[5D10C57]
Một máy bay cất cánh từ sân bay Nội Bài lúc $22$ giờ $45$ phút đêm và hạ cánh tại sân bay Tân Sơn Nhất sau $2$ giờ $10$ phút bay. Hỏi máy bay hạ cánh lúc mấy giờ mấy phút sáng hôm sau? Tính số phút tính từ $0$ giờ của ngày mới.
\shortans{55}
\loigiai{
\textbf{Đáp số:} $55$\\
Thời điểm hạ cánh: $22$ giờ $45$ phút $+ 2$ giờ $10$ phút $= 24$ giờ $55$ phút $= 0$ giờ $55$ phút ngày hôm sau $= 55$ phút.
}
\end{ex}
